%% sections/01_introduction.tex
\section{Introduction}
\label{sec:intro}

Medical licensing examinations such as the \gls{usmle} represent one of the most demanding knowledge challenges in professional education. Candidates must synthesize clinical reasoning across dozens of specialties under timed, high-stakes conditions.Traditional preparation relies on self-directed study, peer quizzing, and expensive tutoring --- resources not equally accessible to all candidates.

This paper presents \gls{emma}, a conversational medical study agent built to close that gap. In explanation mode, students pose clinical questions in natural language and receive responses grounded in passages from eighteen standard medical textbooks. In quiz mode, \gls{emma} presents authentic \gls{usmle}-style questions, evaluates answers, and tracks per-specialty performance over time. A \gls{cf} recommender then steers students toward their weakest areas.

\gls{emma} is built on a five-component \gls{ml} pipeline: a SpaCy \gls{ner} model for clinical entity extraction, a \gls{faiss} vectorstore over 36,723 textbook chunks, a \gls{tfidf} + \gls{lsvc} specialty classifier, a BERTopic clustering model for fine-grained topic routing, and a Surprise-based \gls{cf} recommender. All components feed a locally-run Qwen3-4B-Thinking-2507 \gls{llm}. Figure~\ref{fig:pipeline} illustrates the full architecture.

The central empirical question is whether a textbook-grounded \gls{rag} improves a small local \gls{llm}'s accuracy on MedQA compared to a no-retrieval baseline. A six-run ablation grid varying embedding model, \gls{llm} variant, and \gls{rag} condition shows the answer is condition-dependent: with the right embedding and \gls{llm} pairing, \gls{rag} delivers a $+11$ \gls{pp} accuracy gain; without those conditions, it can hurt. This nuanced finding is the primary contribution of this work.

The remainder is organized as follows: Section~\ref{sec:problem} formalizes the problem. Section~\ref{sec:data} describes data preparation. Sections~\ref{sec:ner} through~\ref{sec:crs} present each pipeline component. Section~\ref{sec:ablation} presents the ablation study, Section~\ref{sec:evaluation} evaluates results, Section~\ref{sec:error} analyzes failure modes, and Section~\ref{sec:conclusion} concludes.

\pagebreak
\begin{strip}
  \centering
  \includegraphics[width=\textwidth, height=0.89\textheight, keepaspectratio]{pipeline}
  \Description{The EMMA five-component pipeline: \gls{ner} extracts clinical entities, \gls{faiss} retrieves textbook chunks, the classifier routes by specialty, BERTopic refines by topic cluster, and the \gls{llm} generates the final response.}
  \captionsetup{hypcap=false}
  \captionof{figure}{EMMA architecture pipeline. A student query passes through \gls{ner}-based query rewriting, \gls{faiss} dense retrieval over 36,723 textbook chunks, specialty classification, BERTopic topic routing, and finally an \gls{llm} that generates a textbook-grounded response. The \gls{cf} recommender operates in parallel on accumulated quiz performance.}
  \label{fig:pipeline}
\end{strip}